\section{Threat Model}

\SCS{} assumes adversaries may:

\begin{enumerate}
  \item submit arbitrary external transactions;
  \item deploy malicious contracts;
  \item call public methods in unexpected orders;
  \item trigger cross-contract calls;
  \item attempt reentrancy;
  \item attempt raw Atomspace mutation;
  \item attempt raw Atomspace enumeration;
  \item attempt to forge caller identity;
  \item attempt to bypass method policies;
  \item attempt to mutate capability registry records;
  \item attempt to exploit nondeterminism;
  \item attempt to leak private state through errors, events, traces, or logs;
  \item attempt to call \PeTTa{}, Prolog, Python, filesystem, process, or network escape hatches;
  \item attempt replay attacks using nonces or signed certificates;
  \item attempt to create duplicate or contradictory balance facts.
\end{enumerate}

\SCS{} assumes the following components are part of the trusted computing base:

\begin{itemize}
  \item contract dispatcher;
  \item method policy registry;
  \item enforcement module;
  \item live capability registry;
  \item guarded storage kernel;
  \item atomic transaction engine;
  \item restricted evaluator implementation;
  \item deterministic serialization and hashing machinery.
\end{itemize}

\SCS{} does not trust ordinary contract code to enforce its own security by convention.
